\section{习题 3.2：二次曲线的类型}
\chr{11}{141—142}{147—148}{304—305}
\K{旋转与平移（在直角坐标中）保持曲线形状，变换后的方程形式决定椭圆、双曲线、抛物线及退化情况。此节第 3—10 题多为证明/分类，原书答案较短，以下给出独立论证。}
\begin{center}
\begin{tikzpicture}[scale=.74]
\begin{scope}[shift={(-4.5,0)}]\draw[->](-2,0)--(2,0);\draw[->](0,-1.8)--(0,1.8);\draw[very thick,accent](0,0)ellipse(1.4 and 1);\node at(0,-2.2){椭圆：两个同号平方项};\end{scope}
\begin{scope}[shift={(0,0)}]\draw[->](-1.9,0)--(1.9,0);\draw[->](0,-1.8)--(0,1.8);\draw[very thick,mint,domain=-1.2:1.2,samples=80] plot({sqrt(1+\x*\x)},\x);\draw[very thick,mint,domain=-1.2:1.2,samples=80] plot({-sqrt(1+\x*\x)},\x);\node at(0,-2.2){双曲线：两平方异号};\end{scope}
\begin{scope}[shift={(4.5,0)}]\draw[->](-1.7,0)--(1.8,0);\draw[->](0,-1.8)--(0,1.8);\draw[very thick,orange!80!black,domain=-1.2:1.2,samples=65]plot({\x*\x},\x);\node at(0,-2.2){抛物线：一平方一次};\end{scope}
\end{tikzpicture}
\end{center}
\readtip{\textbf{退化类型不要漏：}两条相交直线、两条平行直线、重合直线、一个点、空集。判别时先检查坐标系，再看方程的秩与二次型正负性。}

\Q{1}{把一般抛物线写成一个平方项与一个一次项}
\textbf{题意：}$y=4x^2-8x+5$。\hint{配方，只平移不旋转。}
\[
y=4(x-1)^2+1\quad\Longrightarrow\quad x'=x-1,\ y'=y-1.
\]
因此新原点取 $(1,1)$，坐标方向保持不变，即 $I'=[(1,1);(1,0),(0,1)]$，新方程为 $y'=4x'^2$。
\ans{只保留一个平方项和一次项，书后提示的原点 $(1,1)$ 正确。}

\Q{2}{把两个空间二次曲面化为可识别的曲面}
\textbf{(1)} $(2x+y+z)^2-(x-y-z)^2=y-z$。\hint{平方差分解成两个线性因子，再检查残余一次项是否沿零特征方向。}
\[
3x(x+2y+2z)=y-z.
\]
设 $U=\sqrt3x,\;V=\sqrt3(x+2y+2z),\;W=y-z$。三个线性形式独立，方程化为 $UV=W$；进一步设 $u=(U+V)/2,v=(U-V)/2$，即 $u^2-v^2=W$，为马鞍型（双曲抛物面）。
\textbf{(2)} $9x^2-25y^2+16z^2-24xz+80x-60z=0$。\hint{前三个含 $x,z$ 的平方配方得到 $(3x-4z)^2$，另引入独立线性形式作第三坐标。}
\[
(3x-4z)^2-(5y)^2=-(80x-60z).
\]
取 $u=3x-4z,v=5y,w=80x-60z$，三者线性独立；得 $u^2-v^2=-w$，仍为双曲抛物面。
\ans{两小题都是\textbf{马鞍面}；与书后简答一致。}

\Q{3}{不过圆锥顶点的平面截线为何只能是三种圆锥曲线}
\hint{把圆锥取成直角坐标中的标准齐次式 $x^2+y^2=k^2z^2$（$k>0$）；任一不过顶点的平面经适当取轴，可设成 $z=h+mx$，$h\ne0$。}
代入锥面得
\[
(1-k^2m^2)x^2+y^2-2k^2hmx-k^2h^2=0.
\]
二次型的符号由 $1-k^2m^2$ 决定：当 $|km|<1$ 为椭圆（水平截为圆）；$|km|=1$ 退化为一条有实点的抛物线；$|km|>1$ 为双曲线。因 $h\ne0$，这里不会出现两条直线通过顶点的退化情形。
\ans{椭圆、抛物线或双曲线；相应的判别条件已完整给出。}

\Q{4}{单叶双曲面过一条直母线的截平面}
\hint{把截平面当二维参数平面，双曲面方程限制在该平面内是一个二元二次式 $G(u,v)$。}
因为平面已含双曲面的一条直母线，可取这条母线为 $v=0$，故 $G(u,0)\equiv0$，所以 $G(u,v)$ 的 $u^2,u,1$ 项全为零，必可写
\[G(u,v)=v(\alpha u+\beta v+\gamma).
\]
截线是 $v=0$ 与 $\alpha u+\beta v+\gamma=0$ 两条直线（非退化单叶双曲面下两族直母线）；若某种极限情形出现重复因子，应按重合直线计重数。
\ans{证明了过直母线的截平面与单叶双曲面的截线分解。}

\Q{5}{单叶双曲面的所有可能平面截线}
\hint{用标准面 $x^2+y^2-z^2=1$，选择不同方向的截平面即可构造代表。}
水平面 $z=h$：$x^2+y^2=1+h^2$ 是椭圆（圆）；竖直面 $x=h$：$y^2-z^2=1-h^2$，$|h|\neq1$ 时为双曲线，$|h|=1$ 为两条相交直线；斜面 $z=x+h$：$y^2-2hx=1+h^2$，$h\ne0$ 时为抛物线；若 $h=0$ 得两条平行直线 $y=\pm1$。
\ans{\textbf{椭圆、双曲线、抛物线、两条相交直线、两条平行直线}；与原书列举一致。}

\Q{6}{马鞍面可能有哪些平面截线}
\hint{标准形取 $z=x^2-y^2$，直接找适当截面。}
$z=h$ 给出双曲线（$h\ne0$）或两条相交直线（$h=0$）；$y=h$ 给出 $z=x^2-h^2$，是抛物线；$y=x+h$ 给出 $z=-2hx-h^2$，是\textbf{一条直线}（截面内）。
\ans{双曲线、抛物线、两条相交直线或一条直线；与原书一致。}

\Q{7}{二元二次函数的图形随等高线类型变化}
\textbf{题意：}若 $Q(x,y)=a_{11}x^2+a_{22}y^2+2a_{12}xy+2b_1x+2b_2y+c$ 的零等高线为椭圆／双曲线／抛物线，判断 $z=Q(x,y)$。\hint{抛物面类型由 Hessian 的符号与秩决定，不是靠原点常数项判断。}
若零等高线是真椭圆，则 $A$ 定号，$z=Q$ 是椭圆抛物面；若零等高线是真双曲线，则 $A$ 不定号，为双曲抛物面（马鞍面）；若零等高线是真抛物线，则 $A$ 秩 1 且线性项沿零方向非零，经旋转平移可变成 $z=u^2+v$，即斜放的抛物柱面。
\ans{依次为椭圆抛物面、马鞍面、抛物柱面。}

\Q{8}{二次方程的零点恰是 $y$ 轴时的系数限制}
\textbf{题意：}在一般仿射平面坐标系中，$F=0$ 的图像恰是 $x=0$。\hint{先要求 $F(0,y)\equiv0$，再排除出现第二条与 $x=0$ 不同的直线。}
条件 $F(0,y)=a_{22}y^2+2b_2y+c\equiv0$ 导致 $a_{22}=b_2=c=0$；此时
\[F(x,y)=x(a_{11}x+2a_{12}y+2b_1).
\]
若 $a_{12}\ne0$ 或 $b_1\ne0$，第二因子为另一条非 $y$ 轴直线，使零点集合不止 $x=0$。所以 $a_{12}=b_1=0$，并需 $a_{11}\ne0$，最终 $F=a_{11}x^2$。
\ans{只有 $a_{11}$ 非零，其余系数全为零。}

\Q{9}{二次多项式若只表示一条直线，则必是平方}
\hint{平移旋转仿射坐标，把该直线化为 $x=0$，前题说明 $F=kx^2$。}
在原坐标下 $x$ 是一个非零线性形式 $\ell(X)=ax+by+c$，因此
\[F(x,y)=k\ell(x,y)^2,\qquad k\ne0.
\]
把平方系数的模吸入 $a,b,c$，剩余系数只可能为 $+1$ 或 $-1$。
\ans{$\boxed{F=\pm(ax+by+c)^2}$。请把“恰为一条直线”与“分解为两条不同直线”区分开。}

\Q{10}{坐标轴的并集与两相交直线的一般形式}
\subpart{(1) $x$ 轴和 $y$ 轴的并集} 代 $y=0$ 得 $F(x,0)\equiv0$，故 $a_{11}=b_1=c=0$；代 $x=0$ 得 $a_{22}=b_2=0$，因原式确属二次方程须 $a_{12}\ne0$。故 $F=2a_{12}xy$，即 $kxy=0$。
\subpart{(2) 两条相交直线} 作可逆仿射变换，使两直线成为新坐标轴 $u=0,v=0$；由 (1) 得 $F=kuv$。还原得到
\[
\boxed{F=(a_1x+b_1y+c_1)(a_2x+b_2y+c_2),\quad a_1b_2-a_2b_1\ne0.}
\]
\ans{两小题均得证，行列式条件保证两线确实相交而不平行。}

\Q{11}{构造过两条二次曲线全部交点及指定点的新二次曲线}
\textbf{原式：}$F_1=x^2-2y^2+xy+6x-1$，$F_2=2x^2-y^2-x-y$；另过 $P=(2,-2)$。\hint{所有两曲线交点同时满足 $F_1=F_2=0$，故取线性组合 $F=\lambda F_1+\mu F_2$。}
先算 $F_1(2,-2)=3$、$F_2(2,-2)=4$，所以取 $\lambda=4,\mu=-3$，便有 $F(P)=12-12=0$。
\[
4F_1-3F_2=-2x^2-5y^2+4xy+27x+3y-4.
\]
乘 $-1$ 不改零点集合，得到
\[\boxed{2x^2+5y^2-4xy-27x-3y+4=0.}\]
\ans{与原书一致；这种“曲线束”方法对消去所有已知交点尤其高效。}
